% !TeX root = 16_002_Fall24.tex
\footnotesep 12pt
\addtolength{\skip\footins}{3mm}
\parskip 12pt
\parindent 0em
\abovedisplayskip=0pt
\abovedisplayshortskip=0pt

%\newcommand{\jred}[1]{\color{red} #1 \color{black}}
%\newtheorem{lemma}{Lemma}
%\newtheorem{remark}{Remark}
%\newtheorem{examp}{Example}
%\newtheorem{defin}{Definition}
\newtheorem{theo}{Theorem}
%\newtheorem{cor}{Corollary}
%\newtheorem{prop}{Proposition}

\newcommand{\jcode}[1]{\begin{lstlisting} #1 \end{lstlisting}}
\renewcommand\figurename{\sf \Large Fig.}

\makeatletter
\renewcommand\xleftrightarrow[2][]{%
	\ext@arrow 9999{\longleftrightarrowfill@}{#1}{#2}}
\newcommand\longleftrightarrowfill@{%
	\arrowfill@\leftarrow\relbar\rightarrow}
\makeatother
%$F \xleftrightarrow{\text{DFT}} G$
%$F \xleftrightarrow{\text{overlong text}} G$

\DeclareMathOperator{\sinc}{sinc}
\DeclareMathOperator{\rect}{rect}
\newcommand{\jsigma}{u}
\newcommand{\jPi}{\text{rect}}
\newcommand{\jLambda}{\text{tri}}

\newcommand{\js}{\vspace*{.15in}}
\newcommand{\jsm}{\vspace*{-.1in}}
\newcommand{\iLT}[1]{\mathcal{L}^{-1}\left\{ #1 \right\}}
\newcommand{\jLT}[1]{\mathcal{L}\left\{ #1 \right\}}
\newcommand{\jFT}[1]{\mathcal{F}\left\{ #1 \right\}}
\newcommand{\jlongint}{\int_{-\infty}^{\infty}}

\DeclareSymbolFont{cyrletters}{OT2}{wncyr}{m}{n}
\DeclareMathSymbol{\Sha}{\mathalpha}{cyrletters}{"58}

\newcommand{\reals}{{\bf \mathbb R}}
\newcommand{\imags}{{\mathbb C}}
\newcommand{\nonn}{{\mathbb N}}
\newcommand{\diags}{{\mathbb D}}
\newcommand{\symm}{{\mathbb S}}
\newcommand{\ints}{{\mathbb I}}
\newcommand{\expec}{{\mathbb E}}

%\newcommand{\jQ}{\jrww}
\newcommand{\jP}{Q}

\newcommand{\REALS}{{\mathbb R}}
\newcommand{\IMAGS}{{\mathbb C}}
\newcommand{\INTS}{{\mathbb I}}
\newcommand{\tra}{\mbox{\rm tr }}
\newcommand{\wrt}{with respect to}
\newcommand{\dteq}{\stackrel{\scriptstyle \Delta}{\scriptstyle =}}
\newcommand{\mins}{\stackrel{\scriptstyle \min}{\sim}}

\newcommand{\jbm}[1]{\begin{minipage}[t]{#1in} \raggedright}
\newcommand{\jbmd}[1]{\begin{minipage}{#1in}}
\newcommand{\jem}{\end{minipage}}
\newcommand{\be}{\begin{equation}}
\newcommand{\ee}{\end{equation}}
\newcommand{\bjv}{\begin{verbatim}}
\newcommand{\ejv}{\end{verbatim}}
\newcommand{\bth}{\begin{theo}}
%\newcommand{\eth}{\end{theo}}
\newcommand{\bprop}{\begin{prop}}
\newcommand{\eprop}{\end{prop}}
\newcommand{\bcor}{\begin{cor}}
\newcommand{\ecor}{\end{cor}}
\newcommand{\bl}{\begin{lemma}}
\newcommand{\el}{\end{lemma}}
\newcommand{\bre}{\begin{remark} \em }
\newcommand{\ere}{\hspace*{\fill} $\Box$ \end{remark}}
\newcommand{\bex}{\begin{examp}}
\newcommand{\eex}{\end{examp}}
\newcommand{\beex}{\begin{examp}}
\newcommand{\eeex}{\end{examp}}

\newcommand{\jbdef}{\begin{defin}}
\newcommand{\jedef}{\end{defin}}

\newcommand{\bt}{\begin{tabular}}
\newcommand{\et}{\end{tabular}}
\newcommand{\bq}{\begin{quote}}
\newcommand{\eq}{\end{quote}}
\newcommand{\bi}{\begin{itemize}\setlength{\belowdisplayskip}{1pt}}
%\newcommand{\bi}{\begin{itemize}}
\newcommand{\bii}{\begin{itemize}\small}
\newcommand{\ei}{\end{itemize}}
\newcommand{\bee}{\begin{enumerate}}
\newcommand{\eee}{\end{enumerate}}
\newcommand{\bea}{\begin{eqnarray}}
\newcommand{\eea}{\end{eqnarray}}
\newcommand{\beas}{\begin{eqnarray*}}
\newcommand{\eeas}{\end{eqnarray*}}
\newcommand{\ba}{\begin{array}}
\newcommand{\ea}{\end{array}}
\newcommand{\jbas}{\begin{align*}}
\newcommand{\jeas}{\end{align*}}
\newcommand{\bff}{\begin{flushleft}}
\newcommand{\eff}{\end{flushleft}}
\newcommand{\bc}{\begin{center}}
\newcommand{\ec}{\end{center}}
\newcommand{\Hinf}{\mbox{${\cal H}_{\infty}$}~}
\newcommand{\hinf}{\mbox{${\cal H}_{\infty}$}}
\newcommand{\Htwo}{\mbox{${\cal H}_{2}$~}}
\newcommand{\htwo}{\mbox{${\cal H}_{2}$}}
\newcommand{\etal}{{\em et al.}}
\newcommand{\mitt}{Massachusetts Institute of Technology}
\newcommand{\simn}{\stackrel{\min}{\sim}}

\newcommand{\jdiag}{\mbox{\rm diag}}
\newcommand{\jvec}{\mbox{\rm vec}}
\newcommand{\jdk}{$D$--$K$}
%\newcommand{\jom}{{\rm j}\omega}
%\newcommand{\bjom}{({\rm j}\omega)}
\newcommand{\jmu}{$\mu$}
\newcommand{\jsp}{\hspace*{\fill}}
%\newcommand{\smin}{{\scriptstyle -}}
\newcommand{\spos}{{\scriptstyle +}}

\newcommand{\jmatl}[4]{\left[ \ba{r|r} #1 & #2 \\ \hline #3 & #4 \ea \right]}
\newcommand{\jmat}[4]{\left[ \ba{rr} #1 & #2 \\ #3 & #4 \ea \right]}
\newcommand{\jmatc}[4]{\left[ \ba{cc} #1 & #2 \\ #3 & #4 \ea \right]}
\newcommand{\jmatd}[2]{\left[ \ba{ccc} #1 \\ & \ddots \\ && #2 \ea \right]}
\newcommand{\jvecvc}[2]{\left[ \ba{c} #1 \\[.5em] #2 \ea \right]}
\newcommand{\jvechc}[2]{\left[ \ba{cc} #1 & #2 \ea \right]}

\newcommand{\jrxx}{R_{\rm xx}}
\newcommand{\jrzz}{R_{\rm zz}}
\newcommand{\jvxx}{V_{\rm xx}}
\newcommand{\jruu}{R_{\rm uu}}
%\renewcommand{\jrvv}{R_{\rm vv}}
%\renewcommand{\jrww}{R_{\rm ww}}
\newcommand{\jvyy}{V_{\rm yy}}
\newcommand{\jbw}{B_{\rm w}}
\newcommand{\jbu}{B_{\rm u}}
\newcommand{\jcy}{C_{\rm y}}
\newcommand{\jcz}{C_{\rm z}}
\newcommand{\jdyw}{D_{\rm yw}}
\newcommand{\jdzw}{D_{\rm zw}}
\newcommand{\jdyu}{D_{\rm yu}}
\newcommand{\jdzu}{D_{\rm zu}}
\newcommand{\jns}{\textbf{NS}~}
\newcommand{\jnp}{\textbf{NP}~}
\newcommand{\jrs}{\textbf{RS}~}
\newcommand{\jrp}{\textbf{RP}~}
\newcommand{\sgn}{\tt{sgn}}
\newcommand{\jrem}{\noindent \textbf{Remark:}}
\newcommand{\imj}{j}
\newcommand{\alert}[1]{\textbf{\color{blue} #1}}
\newcommand{\alertm}[1]{{\color{blue} \pmb{#1}}}

\newcommand{\jbx}{{\rm x}}
\newcommand{\jby}{{\rm y}}
\newcommand{\jbz}{{\rm z}}
\newcommand{\jfa}{\hspace*{.1in} \forall}

\font\LAsym=cmex10 scaled \magstep2
\def\LARGint{\textfont3=\LAsym\mathchar"1352\nolimits}
%\def\LARGsum{\textfont3=\LAsym\mathchar"1350\nolimits}
\def\LARGsum{\textfont3=\LAsym\mathchar"1350}

\newcommand{\jintI}{\LARGint_{-\infty}^{\infty}}
\newcommand{\jintt}{\LARGint_{0^{-}}^{t}}
\newcommand{\jinti}{\LARGint_{0^{-}}^{\infty}}

\def\labelenumi{\arabic{enumi}.}
\def\labelenumii{(\alph{enumii})}
\def\labelenumiii{(\roman{enumiii})}

\renewcommand{\labelitemiii}{\color{red} $\bullet$\color{black}}
\newcommand{\figref}[1]{Figure~\ref{fig:#1}}
\newcommand{\jw}{j\omega}
\newcommand{\trace}[1]{\mbox{{\rm tr}}\left\{ #1\right\}}
\newcommand{\dw}{\mbox{\rm d}\omega}
\newcommand{\sigmax}[1]{\overline{\sigma}#1}
\newcommand{\reg}{\raisebox{1ex}{{\scriptsize $\bigcirc \!\!\!\!\!$\tiny R}\ }}

\newcommand{\jfig}[4]{
\begin{minipage}[t]{.49\textwidth}
\vspace*{#4}
\begin{center}
\includegraphics[width = 1\textwidth]{#1}
\vspace*{-.25in} \captin[#2]{#2} \label{#3}
\end{center}
\end{minipage}
}

\newcommand{\smax}{\overline{\sigma}}
\newcommand{\smin}{\underline{\sigma}}

\newcommand{\njblue}[1]{{\color{blue}\textbf{#1}}}
\newcommand{\jblue}{\color{blue}\bf}
\newcommand{\njgreen}[1]{\textcolor{orange}{\textbf{#1}}}
\newcommand{\jgreen}{\color{orange}\bf}
\newcommand{\jred}[1]{{\color{red} #1}}
\newcommand{\njred}[1]{{\color{red}\textbf{#1}}}
\newcommand{\njmag}[1]{{\color{purple}\textbf{#1}}}
\newcommand{\njgray}[1]{{\color{lightgray} #1}}
\definecolor{darkbrown}{rgb}{0.4, 0.26, 0.13}
\definecolor{darkgreen}{rgb}{0,0.5,0}
\definecolor{darkred}{rgb}{.86,.08,.24}
\definecolor{gold}{rgb}{1.0, 0.84, 0.0} 
\definecolor{lightblue}{HTML}{DEF2FB}
\definecolor{darkblue}{HTML}{BCE5FB}
\definecolor{lightaqua}{HTML}{defbf6}
\definecolor{darkaqua}{HTML}{bcfbf2}
\definecolor{lightbluetext}{HTML}{E6F7FF}
\definecolor{darkbluetext}{HTML}{D3EFFE}
\definecolor{bluetext}{HTML}{2E73A3}
\definecolor{grey}{HTML}{AFAFAF}


\vfuzz8pt % Don't report over-full v-boxes if over-edge is small
\hfuzz8pt % Don't report over-full h-boxes if over-edge is small

\renewcommand{\jmatc}[4]{\left[ \ba{cc} #1 & #2 \\[.25em] #3 & #4 \ea \right]}%
\newcommand{\jmatcl}[4]{\left[ \ba{c|c} #1 & #2 \\\hline #3 & #4 \ea \right]}%
\newcommand{\jmatr}[4]{\left[ \ba{rr} #1 & #2 \\[.25em] #3 & #4 \ea \right]} %
\newcommand{\jmthr}[9]{\left[ \ba{rrr} #1 & #2 & #3 \\ #4 & #5 & #6 \\ #7 & #8 & #9\ea \right]}
\newcommand{\jvthr}[3]{\left[ \ba{r} #1 \\ #2 \\ #3\ea \right]} %
\newcommand{\jhthr}[3]{\left[ \ba{rrr} #1 & #2 & #3\ea \right]} %
\newcommand{\jvecc}[2]{\left[ \ba{cc} #1 & #2 \ea \right]} %
\newcommand{\jcolc}[2]{\left[ \ba{c} #1 \\ #2 \ea \right]} %
\newcommand{\jcolr}[2]{\left[ \ba{r} #1 \\ #2 \ea \right]} %
\newcommand{\jrowr}[2]{\left[ \ba{rr} #1 & #2 \ea \right]} %
\newcommand{\jcoltc}[3]{\left[ \ba{c} #1 \\ #2 \\ #3\ea \right]} %

%define \def\mathbi#1{\mathbf{\em #1}}
\newcommand{\XX}{{\bf XX}}
\newcommand{\jmo}{{\cal M}_0}
\newcommand{\jmc}{{\cal M}_c}
\newcommand{\jq}{\frac{1}{2}\rho U_0^2}
\newcommand{\jh}{\frac{1}{2}}
\newcommand{\jar}{{\color{blue}$\mathbf{\Rightarrow}$}~}
\newcommand{\njar}{{\color{blue}\mathbf{\Rightarrow}}~}
\newcommand{\jxstar}{\njbx^\star}
\newcommand{\xstar}{\mathbf{x}^\star}
\newcommand{\hxstar}{\mathbf{\hat x}^\star}
\newcommand{\njbI}{\mathbf{I}}
\newcommand{\njbJ}{\mathbf{J}}
\newcommand{\njbA}{\mathbf{A}}
\newcommand{\njbB}{\mathbf{B}}
\newcommand{\njbC}{\mathbf{C}}
\newcommand{\njbD}{\mathbf{D}}
\newcommand{\njbF}{\mathbf{F}}
\newcommand{\njbH}{\mathbf{H}}
\newcommand{\njbxh}{\widehat{\mathbf{x}}}
\newcommand{\njbxt}{\tilde{\mathbf{x}}}
\newcommand{\njbxhn}{\hat{\mathbf{x}}}
\newcommand{\njbyh}{\widehat{\mathbf{y}}}
\newcommand{\njbu}{\mathbf{u}}
\newcommand{\njbf}{\mathbf{f}}
\newcommand{\njbg}{\mathbf{g}}
\newcommand{\njbh}{\mathbf{h}}
\newcommand{\njba}{\mathbf{a}}
\newcommand{\njbb}{\mathbf{b}}
\newcommand{\njbi}{\mathbf{i}}
\newcommand{\njbe}{\mathbf{e}}
\newcommand{\njbd}{\mathbf{d}}
\newcommand{\njbp}{\mathbf{p}}
\newcommand{\njbq}{\mathbf{q}}
\newcommand{\njbr}{\mathbf{r}}
\newcommand{\njby}{\mathbf{y}}
\newcommand{\njbv}{\mathbf{v}}
\newcommand{\njbw}{\mathbf{w}}
\newcommand{\njbx}{\mathbf{x}}
\newcommand{\njbX}{\mathbf{X}}
\newcommand{\njbY}{\mathbf{Y}}
\newcommand{\njbz}{\mathbf{z}}
\newcommand{\njbs}{\mathbf{s}}
\newcommand{\njbm}{\mathbf{m}}
\newcommand{\njbn}{\mathbf{n}}
\newcommand{\njbc}{\mathbf{C}}
\newcommand{\ndegs}{$^\circ$~}
\newcommand{\degs}{^\circ}
\newcommand{\njbnu}{\boldsymbol{\nu}}
\newcommand{\jfbox}[1]{\fbox{\begin{Beqnarray*} #1 \end{Beqnarray*}}}
\newcommand{\jfboxns}[1]{\fbox{\begin{Beqnarray} #1 \end{Beqnarray}}}
\newcommand{\jobox}[1]{\Ovalbox{\begin{Beqnarray*} #1 \end{Beqnarray*}}}

%\newcommand{\jrxx}{R_{\rm xx}}
%\newcommand{\jruu}{R_{\rm uu}}
\newcommand{\jrww}{R_{\rm ww}}
\newcommand{\jrvv}{R_{\rm vv}}

\newcommand{\nbas}{\begin{align*}}
\newcommand{\neas}{\end{align*}}
\newcommand{\vf}{\vfill}
\newcommand{\jmin}[1]{\bi \item #1 \ei}
\newcommand{\jom}{\mathbf{j}\omega}
\newcommand{\bjom}{(\mathbf{j}\omega)}
\newcommand{\jomp}{\mathbf{j}\omega_{\pi}}
\newcommand{\bjomc}{(\mathbf{j}\omega_{c})}
\newcommand{\bjomp}{(\mathbf{j}\omega_{\pi})}
\newcommand{\jomc}{\mathbf{j}\omega_{c}}
\newcommand{\jtem}{\vfill\item}
\newcommand{\weird}[2]{~\vphantom{#1}^{#2} #1}

\newcommand{\jnl}{\newline}
\newcommand{\jei}{{\em i.e.,~}}
\newcommand{\jd}{\displaystyle}
\def\Tiny{\fontsize{6pt}{6pt}\selectfont}
\def\RTiny{\fontsize{.1pt}{.1pt}\selectfont}

\newcommand{\fo}{\forall~~\omega}
%\newcommand{\degs}{$^\circ$}
\newcommand{\ra}{$\Rightarrow$~}
\newcommand{\dsd}[2]{\frac{\partial #1}{\partial #2}}
\newcommand{\qs}{\frac{1}{2}\rho V_T^2 S}
\newcommand{\csd}[2]{C_{#1_#2}}
\newcommand{\Ma}{\mathbf{M}}
\newcommand{\jl}{\left(}
\newcommand{\jr}{\right)_0}
\newcommand{\jbt}[1]{\begin{tabular}{#1}}
\newcommand{\jet}[1]{\end{tabular}}
\newcommand{\jgtd}{G_{\theta\delta_e}}
\newcommand{\jgqd}{G_{q\delta_e}}
\newcommand{\jhalf}{\frac{1}{2}}
\newcommand{\alertr}[1]{\textbf{\color{red} #1}}

%\newcommand{\ojminus}{\!\!\stackrel{\scriptstyle \ominus}{\vphantom{.}}}
%\newcommand{\ojplus}{\!\!\stackrel{\scriptstyle \oplus}{\vphantom{.}}}
\newcommand{\njminus}{_{k|k-1}}
\newcommand{\sjminus}{_{k+1|k}}
\newcommand{\njplus}{_{k|k}}
\newcommand{\ojplus}{_{k-1|k-1}}
\newcommand{\jrect}{\text{\tt rect}}
\newcommand{\jtri}{\text{\tt tri}}

\renewcommand{\XX}[1]{{\bf \color{red} XX #1 XX}}
\newcommand{\jbox}[2]{{\color{#2} \framebox{\color{#2} #1}}}

\newcommand\Tstrut{\rule{0pt}{2.6ex}}         % = `top' strut
\newcommand\Bstrut{\rule[-0.9ex]{0pt}{0pt}}   % = `bottom' strut


\usepackage{xifthen}% provides \isempty test
\newcommand{\optarg}[1][]{%
  \ifthenelse{\isempty{#1}}%
    {Cont'd}% if #1 is empty
    {#1}% if #1 is not empty
}

%% Frames
\newcommand{\jfr}[2]
{\subsection{\optarg{#1}} \begin{frame} #2 \end{frame}}

\newcommand{\jfrt}[2]{\subsection{#1} \begin{frame}{#1}
    \scriptsize #2 \end{frame}}
\newcommand{\jfrnt}[2]{\subsection{~~~~Cont'd} \begin{frame}{}
    \scriptsize #2 \end{frame}} 

\newcommand{\jfrh}[2]
{\subsection{\optarg{#1}} \begin{frame}<handout:0>{#1} \scriptsize #2 \end{frame}}

\newcommand{\jfrf}[2]
{\subsection{\optarg{#1}} \begin{frame}[fragile]{#1} \scriptsize #2 \end{frame}}

\newcommand{\jfrb}[2]
{\subsection{\optarg{#1}} \begin{frame}{#1} \bi #2 \ei \end{frame}}

\newcommand{\jcols}[4]
{\begin{columns}[b]
\begin{column}{#3\textwidth} #1 \end{column}
\begin{column}{#4\textwidth} #2 \end{column}
\end{columns}}

\newcommand{\jcolst}[4]
{\begin{columns}[T]
\begin{column}{#3\textwidth} #1 \end{column}
\begin{column}{#4\textwidth} #2 \end{column}
\end{columns}}

\newcommand{\jcolsb}[4]
{\begin{columns}[b] \begin{column}{#3\textwidth} #1 \vspace{0pt} \end{column}
\begin{column}{#4\textwidth} #2 \vspace{0pt} \end{column} \end{columns}}


%% Thms
\theoremstyle{definition}
\newtheoremstyle{mythmstyle}{}{}{}{}{\bf}{}{ }
{\thmname{#1}\thmnumber{ #2}%
	\thmnote{: #3\addcontentsline{toc}{subsubsection}{{#1} - #3}}: }
\theoremstyle{mythmstyle}
\newtheorem{example}{Example}
\newtheorem{theorem}{Theorem}

\newcommand{\jsec}[1]{\section{#1}} 
\newcounter{mytheorem}[section]
\def\themytheorem{\Lect.\arabic{mytheorem}}

\renewcommand{\jmatc}[4]{\left[ \ba{cc} #1 & #2 \\[.25em] #3 & #4 \ea \right]}%
\newcommand{\jmatcs}[4]{\left[ \ba{cc} #1 & #2 \\[1em] #3 & #4 \ea \right]}

\newcommand{\jreset}[1]{\setcounter{lect}{#1}\setcounter{section}{0}\setcounter{figure}{0}\setcounter{equation}{0}\setcounter{example}{0}}

\newcommand{\jreading}[1]{
	\vfill
	\subsection*{Recommended reading}
	A.\ Oppenheim and A.\ Willsky: #1 
	\vfill
}

\pgfmathdeclarefunction{H}{2}{%
	\pgfmathparse{ (#1 > #2) }%
}

\pgfkeys{/pgf/declare function={% heaviside(var,left limit,right limit. amplitude)
		heaviside(\x,\l,\r,\a)=max( 0 , ( \x>= \l ?(\x> \r?0:\a):0);
	} 
}

\makeatletter
%\pgfmathdeclarefunction{p}{1}{\edef\pgfmathresult{\ifdim#1pt<\z@0\else\ifdim#1pt>1pt 0\else1\fi\fi}}
\pgfmathdeclarefunction{p}{1}{\edef\pgfmathresult{\ifdim#1pt<\z@0\else1\fi}}
\makeatother

\newcommand{\sneakylabel}[1]{
	\coordinate (sneakylabel{#1}) at (current plot begin);
	\pgfplotsset{
		after end axis/.append code={
			\node [anchor=east] at (sneakylabel{#1}){#1};
		}
	}
}

\newsavebox{\mybottombox} % Box to save the text of the command 
\newlength{\mybottomlength} % The length of our text inside the command
\newlength{\availafter} % The available length left on the page after placing our text

% Optional argument is the minimum length after the nobottom text for not pagebreak. Change it to your needs
\newcommand{\nobottom}[2][60pt]{\savebox{\mybottombox}{\vbox{#2}}\setlength{\mybottomlength}{\ht\mybottombox}%
	\setlength{\availafter}{\dimexpr\textheight-\mybottomlength-\pagetotal\relax}\ifdim\availafter<#1%
	\pagebreak\noindent\usebox{\mybottombox}%
	\else%
	\noindent\usebox{\mybottombox}%
	\fi%
}%

%% Figures
\usepackage{animate}
\newcommand{\reserveandshow}[2][]{%
    \phantom{\includegraphics<+>[#1]{#2}}%
    \includegraphics<+->[#1]{#2}%
}
\newcommand{\reserveandshown}[4][]{%
    \phantom{\includegraphics<1-#3>[#1]{#2}}%
    \includegraphics<#4->[#1]{#2}%
}

\newcommand{\inputTikZ}[2]{%  
     \scalebox{#1}{\input{#2}}  
}


%% Boxes
\newcommand{\eqbox}[1]{\begin{empheq}[box={\mymath[colback=red!30,drop lifted shadow, sharp corners]}]{equation}
		\jd #1 \end{empheq}}

\newcommand{\teqbox}[1]{
\eqbox{\notag \text{\begin{minipage}{5in} \begin{center} #1 \end{center} \end{minipage}}}
}

%% TeX-side calculations and formatted numeric output
\usepackage{xfp}
\usepackage{siunitx}
\sisetup{
  round-mode=places,
  round-precision=3
}
\newcommand{\calc}[1]{\fpeval{#1}}
\newcommand{\printnum}[2][2]{%
  \num[round-precision=#1]{\fpeval{#2}}%
}
% Signed printer: prints +x or -x, avoiding "+ -" artifacts.
\newcommand{\signednum}[2][2]{%
  \ifnum\fpeval{(#2) < 0}=1
    -\printnum[#1]{abs(#2)}%
  \else
    +\printnum[#1]{#2}%
  \fi
}
\newcommand{\printadd}[2]{%
  \pgfmathsetmacro{\temp}{#1 + #2}%
  \pgfmathprintnumber[fixed, precision=2]{\temp}%
}

\newcommand{\printsub}[2]{%
  \pgfmathsetmacro{\temp}{#1 - #2}%
  \pgfmathprintnumber[fixed, precision=2]{\temp}%
}

\newcommand{\printmul}[2]{%
  \pgfmathsetmacro{\temp}{#1 * #2}%
  \pgfmathprintnumber[fixed, precision=2]{\temp}%
}

\newcommand{\printdiv}[2]{%
  \pgfmathsetmacro{\temp}{#1 / #2}%
  \pgfmathprintnumber[fixed, precision=2]{\temp}%
}

\newcommand{\printtan}[1]{%
  \pgfmathparse{tan(#1)}%
  \pgfmathprintnumber[fixed,precision=3]{\pgfmathresult}%
}

\newcommand{\printatan}[1]{%
  \pgfmathparse{atan(#1)}%
  \pgfmathprintnumber[fixed,precision=3]{\pgfmathresult}%
}

\newcommand{\printsqrt}[1]{%
  \pgfmathparse{sqrt(#1)}%
  \pgfmathprintnumber[fixed,precision=3]{\pgfmathresult}%
}

\newcommand{\printlog}[1]{%
  \pgfmathparse{log10(#1)}%
  \pgfmathprintnumber[fixed,precision=3]{\pgfmathresult}%
}

\newcommand{\trim}[1]{\pgfmathprintnumber[fixed, precision=2]{#1}}

\newcommand{\findroots}[3]{%
  \edef\posroot{%
    \fpeval{(-(#2) + sqrt((#2)^2 - 4*(#1)*(#3)))/(2*(#1))}%
  }%
  \edef\negroot{%
    \fpeval{(-(#2) - sqrt((#2)^2 - 4*(#1)*(#3)))/(2*(#1))}%
  }%
}

\allowdisplaybreaks
